\section{Introduction}
近年来，大语言模型（LLMs）正从“文本生成器”走向“决策主体”，被用于推荐、定价、谈判与政策咨询等存在利益博弈和反馈循环的场景。在这类战略环境中，模型不仅要响应局部任务信息，还要理解行动如何经由经济机制影响他人，并反馈到自身收益与风险。数据市场中的隐私决策是一个天然测试床：个体收益依赖多主体互动，平台又通过数据推断与个性化决策，使“信息—价格—福利”形成跨阶段因果链与隐蔽外部性。

近期研究已从博弈论和经济学视角考察 LLM 的隐私决策（如用户是否共享、平台如何定价、系统利润与福利如何变化），但主流评估仍偏向终局指标，例如分享率、利润或社会福利是否提升。问题在于：当行为“看起来正确”时，我们难以区分模型是理解了外部性机制，还是依赖语言启发式、语料模式匹配或提示诱导“碰巧做对”。在隐私外部性场景下，这种不可区分性尤其危险，因为外部性常常隐蔽、反直觉且对扰动和交互路径敏感，终局正确可能掩盖机制偏差，并在环境变化时导致失稳甚至市场失败。换言之，“结果正确”不等于“机制正确”。

为弥补这一缺口，本文提出一个机制导向的隐私外部性 LLM benchmark。核心立场是：评估中心应从“输出像不像人/结果好不好看”转向“机制是否被理解”。据此，我们构建三个难度递进的可计算经济场景：场景一关注推断外部性（inference externality）下的数据定价与过度分享（主案例）；场景二关注价格传导型外部性下的双侧互动与均衡一致性；场景三引入中介与机制设计，考察复杂制度下多主体系统的动态稳定性与失效边界。每个场景都配备可复现的理论求解器，生成可验证的 ground truth（均衡策略、价格与福利分解等），避免依赖人工标注与主观偏好。

在评测协议上，我们将 LLM 以自然语言交互方式嵌入博弈过程，使其扮演消费者、平台或中介等策略主体，并在统一管线下比较其行为与理论解。除终局偏差（outcome deviation）外，我们进一步引入两类诊断维度：策略结构一致性（structural consistency，如预测分享集合与理论均衡集合的一致性）与动态过程性质（process properties，如收敛速度、稳定性、路径依赖），用于识别“高分但机制错”的伪理解，并评估模型是否能在重复交互中形成稳定策略。

此外，我们在主案例中设计“提示词信息密度梯度”实验（v0–v5）：在经济环境与理论解不变的前提下，逐步从基础规则说明提升到机制概念注入、结构化解释与公式框架，观察模型行为随信息结构变化的响应曲线，从而诊断其是机制推断还是文本诱导下的表面正确。

我们在 GPT、DeepSeek、Qwen 三个家族的 12 个前沿模型上进行了系统评测。结果显示：仅增加参数与细节并不会稳定提升机制一致性，反而可能诱发“结果合理但结构偏离”的伪理解；引入“推断外部性”这一关键机制概念后，模型与理论解的一致性出现显著跃迁。对部分弱模型引入多轮虚拟博弈（fictitious play）后，其机制学习能力、动态收敛性与均衡一致性均有改善。整体上，机制导向的诊断评测能揭示传统终局指标掩盖的差异，更准确刻画模型在机制理解、鲁棒性与动态稳定性上的能力边界。

综上，本文贡献包括：第一，提出机制导向的隐私外部性 LLM benchmark，并构建三类可计算评测场景；第二，为各场景提供可复现理论求解器与可验证 ground truth，提升评测客观性与可复现性；第三，建立统一的多主体自然语言交互协议与诊断指标体系，显式区分“终局正确”与“机制正确”，并纳入动态收敛与稳定性评估；第四，在 12 个前沿模型上给出关于“参数≠理解”、机制概念注入拐点及虚拟博弈学习效果的实证结论，为后续研究与模型改进提供可解释基准。
